% Chapter 12: Complex Analysis with Euler's Formula
% ch12_euler_formula.tex

\chapter{Complex Analysis: Euler's Formula Extended}

\begin{historical}
Euler's formula $e^{ix} = \cos x + i \sin x$ opened the door to complex analysis.
This chapter extends the formula to complex exponents, logarithms, and powers,
forming the foundation of modern complex function theory that Riemann, Cauchy, and
Weierstrass would develop further in the 19th century.
\end{historical}

\section{Complex Exponentiation}

Euler's formula extends naturally to complex exponents. For $z = a + ib$:
\[
e^{a+ib} = e^a e^{ib} = e^a (\cos b + i \sin b)
\]

The modulus and argument are:
\[
|e^z| = e^a, \quad \arg(e^z) = b \pmod{2\pi}
\]

\subsection{Properties of Complex Exponential}

\begin{theorem}
The complex exponential function satisfies:
\begin{enumerate}
\item $e^{z_1 + z_2} = e^{z_1} \cdot e^{z_2}$ for all $z_1, z_2 \in \mathbb{C}$
\item $e^0 = 1$
\item $e^z \neq 0$ for all $z \in \mathbb{C}$
\item $e^{z + 2\pi i} = e^z$ (periodicity with period $2\pi i$)
\end{enumerate}
\end{theorem}

\begin{proof}
Property 1 follows from the series definition or from the addition formula for exponentials.
Property 2 is immediate. Property 3 follows since $|e^z| = e^a > 0$.
Property 4 follows from Euler's formula: $e^{2\pi i} = \cos(2\pi) + i\sin(2\pi) = 1$.
\end{proof}

\subsection{Taylor Series for Complex Exponential}

\[
e^z = \sum_{n=0}^{\infty} \frac{z^n}{n!} = 1 + z + \frac{z^2}{2!} + \frac{z^3}{3!} + \cdots
\]

This series converges absolutely for all $z \in \mathbb{C}$.

\begin{lstlisting}[caption={Complex exponential}, label=lst:complex_exp]
from euler import complex_exponential, complex_logarithm, complex_power
import math

z = 1 + 2j
print(f"Complex exponential: e^{z} = {complex_exponential(z)}")
print(f"  |e^z| = {abs(complex_exponential(z)):.6f}")

w = 3 + 4j
log_w = complex_logarithm(w)
print(f"Logarithm: log({w}) = {log_w}")

base = 2 + 3j
exp = 1 + 1j
print(f"Complex power: ({base})^({exp}) = {complex_power(base, exp)}")
\end{lstlisting}

\section{The Complex Logarithm}

The complex logarithm is multi-valued:
\[
\log z = \ln|z| + i(\arg z + 2\pi k), \quad k \in \mathbb{Z}
\]

The \textbf{principal branch} restricts $\arg z \in (-\pi, \pi]$, denoted $\operatorname{Log} z$.
This makes $\log$ single-valued but discontinuous along the negative real axis (the branch cut).

\begin{theorem}[Complex Logarithm]
For $z \neq 0$:
\[
\operatorname{Log} z = \ln|z| + i \operatorname{Arg}(z)
\]
where $\operatorname{Arg}(z) \in (-\pi, \pi]$ is the principal argument.
\end{theorem}

\subsection{Branch Cuts}

\begin{itemize}
\item \textbf{Principal branch cut:} Negative real axis $(-\infty, 0]$
\item \textbf{Other branches:} Can choose different cuts, e.g., positive real axis
\item \textbf{Branch points:} $z = 0$ and $z = \infty$ are branch points of $\log z$
\end{itemize}

\begin{lstlisting}[caption={Complex logarithm}, label=lst:complex_log]
from euler import complex_logarithm
import math, cmath

# Principal branch
for z in [1, -1, 1j, -1j, 1+1j, -1+1j, -1-1j, 1-1j]:
    log_z = complex_logarithm(z)
    expected = cmath.log(z)
    print(f"log({z:6.2f}) = {log_z.real:.4f} + {log_z.imag:.4f}i")

# Verify: exp(Log(z)) = z
print("\nVerification: e^(Log(z)) = z")
for z in [1, 2+3j, -1+1j]:
    result = complex_exponential(complex_logarithm(z))
    print(f"  e^({z}) = {result.real:.6f} + {result.imag:.6f}i")
\end{lstlisting}

\subsection{Logarithm Properties}

\begin{theorem}
For the principal branch:
\begin{enumerate}
\item $\operatorname{Log}(z_1 z_2) = \operatorname{Log} z_1 + \operatorname{Log} z_2 + 2\pi i k$ for some $k \in \{-1, 0, 1\}$
\item $\operatorname{Log}(z^n) = n \operatorname{Log} z + 2\pi i k$ for $k \in \{-n+1, \ldots, n-1\}$
\item $\frac{d}{dz} \operatorname{Log} z = \frac{1}{z}$ for $z \notin (-\infty, 0]$
\end{enumerate}
\end{theorem}

\begin{example}
$\operatorname{Log}((-1)(-1)) = \operatorname{Log}(1) = 0$, but $\operatorname{Log}(-1) + \operatorname{Log}(-1) = i\pi + i\pi = 2\pi i$.
The difference is $2\pi i$, showing the logarithm is not a true homomorphism.
\end{example}

\section{Complex Powers}

For complex $z$ and $w$, define:
\[
z^w = e^{w \log z}
\]

Using the principal branch:
\[
z^w = \exp(w (\ln|z| + i \operatorname{Arg} z))
\]

\begin{theorem}[Complex Power]
For $z \neq 0$ and any complex $w$:
\[
z^w = e^{w \operatorname{Log} z} = \exp(w(\ln|z| + i \operatorname{Arg} z))
\]
\end{theorem}

\subsection{Multivalued Nature of Complex Powers}

Since $\log z$ is multi-valued, $z^w$ is also multi-valued:
\[
z^w = \exp(w(\ln|z| + i(\operatorname{Arg} z + 2\pi k))), \quad k \in \mathbb{Z}
\]

\begin{example}[Special Cases]
\begin{itemize}
\item $i^i = e^{i \cdot i\pi/2} = e^{-\pi/2} \approx 0.2079$ (principal value)
\item $(-1)^{1/2} = \pm i$
\item $0^0$ is indeterminate in complex analysis
\end{itemize}
\end{example}

\section{Euler's Formula for Hyperbolic Functions}

Euler's formula connects to hyperbolic functions:
\[
\cosh x = \frac{e^x + e^{-x}}{2}, \quad \sinh x = \frac{e^x - e^{-x}}{2}
\]

Setting $x \to ix$:
\[
\cos x = \cosh(ix), \quad \sin x = -i \sinh(ix)
\]

\subsection{Euler's Formulas for Hyperbolic Functions}

\begin{theorem}
For any complex $z$:
\[
e^z = \cosh z + \sinh z, \quad e^{-z} = \cosh z - \sinh z
\]
\[
\cosh z = \frac{e^z + e^{-z}}{2}, \quad \sinh z = \frac{e^z - e^{-z}}{2}
\]
\end{theorem}

\begin{lstlisting}[caption={Hyperbolic identities}, label=lst:hyperbolic]
import math, cmath

# Verify: cos(ix) = cosh(x)
for x in [0.5, 1.0, 2.0]:
    cos_ix = cmath.cos(1j * x)
    cosh_x = math.cosh(x)
    print(f"cos(i{x:.1f}) = {cos_ix.real:.6f} + {cos_ix.imag:.6f}i  (cosh({x:.1f}) = {cosh_x:.6f})")
\end{lstlisting}

\subsection{Trigonometric-Hyperbolic Identities}

\begin{theorem}
\begin{align*}
\cos(ix) &= \cosh x \\
\sin(ix) &= i\sinh x \\
\cos(x + iy) &= \cos x \cosh y - i\sin x \sinh y \\
\sin(x + iy) &= \sin x \cosh y + i\cos x \sinh y
\end{align*}
\end{theorem}

\section{Branch Cuts and Multi-Valued Functions}

The complex logarithm and power functions are multi-valued. When working with them,
we must choose a branch. The principal branch cuts along the negative real axis
are the standard convention.

\begin{itemize}
\item \textbf{Logarithm branch cut:} Negative real axis, $(-\infty, 0]$
\item \textbf{Power branch cut:} Depends on the branch of the logarithm used
\item \textbf{Continuity:} Functions are continuous on $\mathbb{C} \setminus \text{branch cut}$
\end{itemize}

\subsection{Riemann Surfaces}

\begin{definition}
A \textbf{Riemann surface} for a multi-valued function is a connected complex manifold
on which the function becomes single-valued and holomorphic.
\end{definition}

\begin{example}
The Riemann surface for $\log z$ is an infinite spiral covering the punctured plane $\mathbb{C} \setminus \{0\}$.
Each turn around the origin ascends to a different sheet.
\end{example}

\section{Euler's Formula and Differential Equations}

Euler's formula solves the differential equation $y' = iy$ with $y(0) = 1$:
\[
y(x) = e^{ix} = \cos x + i \sin x
\]

More generally, $y' = \lambda y$ has solution $y = e^{\lambda x}$.
For $\lambda = a + ib$: $y = e^{ax}(\cos bx + i \sin bx)$.

\subsection{Linear ODEs with Complex Coefficients}

\begin{theorem}
The general solution of $y'' + \omega^2 y = 0$ is:
\[
y(t) = A e^{i\omega t} + B e^{-i\omega t} = C_1 \cos(\omega t) + C_2 \sin(\omega t)
\]
\end{theorem}

\section{Analytic Functions}

\subsection{Cauchy-Riemann Equations}

\begin{theorem}[Cauchy-Riemann Equations]
A function $f(z) = u(x,y) + iv(x,y)$ is holomorphic at $z_0$ if and only if:
\[
\frac{\partial u}{\partial x} = \frac{\partial v}{\partial y}, \quad \frac{\partial u}{\partial y} = -\frac{\partial v}{\partial x}
\]
at $z_0$, and the partial derivatives are continuous.
\end{theorem}

\subsection{Holomorphic Functions}

\begin{definition}
A function is \textbf{holomorphic} on an open set $U$ if it is complex differentiable at every point in $U$.
\end{definition}

\begin{theorem}
The following functions are holomorphic on their domains:
\begin{itemize}
\item Polynomials: $\mathbb{C} \to \mathbb{C}$
\item Rational functions: $\mathbb{C} \setminus \{\text{poles}\}$
\item Exponential: $e^z$: $\mathbb{C} \to \mathbb{C} \setminus \{0\}$
\item Logarithm: $\operatorname{Log} z$: $\mathbb{C} \setminus (-\infty, 0]$
\end{itemize}
\end{theorem}

\section{Applications in Signal Processing}

Euler's formula is fundamental to Fourier analysis:
\[
f(t) = \sum_{n=-\infty}^{\infty} c_n e^{in\omega_0 t}
\]

Each term $c_n e^{in\omega_0 t}$ represents a rotating vector (phasor) in the complex plane.
This representation is essential in:
\begin{itemize}
\item \textbf{Digital signal processing:} FFT algorithms decompose signals into complex exponentials
\item \textbf{Communications:} Modulation schemes (AM, FM, QAM) use complex carriers
\item \textbf{Control theory:} Transfer functions use $e^{st}$ for stability analysis
\end{itemize}

\subsection{Laplace Transform}

The Laplace transform uses $e^{-st}$ with $s = \sigma + i\omega$:
\[
\mathcal{L}\{f(t)\} = F(s) = \int_0^{\infty} f(t) e^{-st} dt
\]

\section{Computational Complexity}

\subsection{Complex Exponential}
Computing $e^{a+ib}$:
\begin{itemize}
\item Time: $O(1)$ for fixed precision (using CORDIC or Taylor series with reduction)
\item Memory: $O(1)$
\item Accuracy: Requires careful argument reduction for large $|b|$
\end{itemize}

\subsection{Complex Logarithm}
Computing $\operatorname{Log} z$:
\begin{itemize}
\item Time: $O(1)$ (involves one sqrt, one log, and one atan2)
\item Memory: $O(1)$
\item Special case: Branch cut requires careful handling for negative reals
\end{itemize}

\subsection{Complex Power}
Computing $z^w = e^{w \operatorname{Log} z}$:
\begin{itemize}
\item Time: $O(1)$ after computing logarithm
\item Memory: $O(1)$
\end{itemize}

\section{Frequently Asked Questions}

\begin{itemize}
\item \textbf{Q: What is the principal branch of the complex logarithm?} \\
The complex logarithm is multi-valued: $\log z = \ln|z| + i(\arg z + 2\pi k)$. The principal branch restricts $\arg z \in (-\pi, \pi]$, denoted $\operatorname{Log} z$. This makes $\log$ single-valued but discontinuous along the negative real axis (the branch cut).

\item \textbf{Q: How do I compute $z^w$ for complex numbers?} \\
Define $z^w = e^{w \log z}$. Using the principal branch: $z^w = \exp(w (\ln|z| + i \operatorname{Arg} z))$. Note that complex exponentiation is not commutative: $z^w \neq w^z$ in general, and $(z^a)^b \neq z^{ab}$ due to branch cuts.

\item \textbf{Q: Why does $e^{i\pi} = -1$ and not some other value?} \\
From Euler's formula: $e^{i\pi} = \cos \pi + i \sin \pi = -1 + 0i = -1$. The cosine and sine values at $\pi$ are determined by the Taylor series definitions, which are consistent with the geometric unit circle interpretation.

\item \textbf{Q: What is the modulus of $e^z$ for complex $z = a + ib$?} \\
$|e^{a+ib}| = |e^a e^{ib}| = e^a |e^{ib}| = e^a \cdot 1 = e^a$. The imaginary part contributes only to the phase (angle), not the magnitude. This is why $e^{ix}$ always lies on the unit circle.

\item \textbf{Q: How does Euler's formula relate to the exponential map in group theory?} \\
The map $x \mapsto e^{ix}$ is a group homomorphism from $(\mathbb{R}, +)$ to the circle group $S^1$ (complex numbers of modulus 1). Its kernel is $2\pi\mathbb{Z}$, giving the isomorphism $\mathbb{R}/2\pi\mathbb{Z} \cong S^1$.

\item \textbf{Q: Why is the complex logarithm multi-valued?} \\
Because $e^z$ is periodic with period $2\pi i$: $e^{z+2\pi i k} = e^z$ for any integer $k$. Therefore, the inverse function must account for all these preimages, giving $\log z = \ln|z| + i(\arg z + 2\pi k)$.

\item \textbf{Q: What is the difference between holomorphic and analytic functions?} \\
In complex analysis, these are equivalent: a function is holomorphic if and only if it is analytic (can be represented by a convergent power series). This is a fundamental difference from real analysis.
\end{itemize}

\section{Exercises}

\begin{exercise}
Prove that $|e^{a+ib}| = e^a$ for all real $a, b$.
\end{exercise}

\begin{exercise}
Show that $\operatorname{Log}(z_1 z_2) \neq \operatorname{Log} z_1 + \operatorname{Log} z_2$ in general by finding a counterexample.
\end{exercise}

\begin{exercise}
Compute $i^i$ using the principal branch and verify it is real.
\end{exercise}

\begin{exercise}[Proof]
Prove the hyperbolic identities $\cos(ix) = \cosh x$ and $\sin(ix) = i\sinh x$.
\end{exercise}

\begin{exercise}
Solve the differential equation $y' = (1+i)y$ with $y(0) = 1$ using Euler's formula.
\end{exercise}

\begin{exercise}[Computational]
Implement the complex logarithm using the principal branch and verify that $e^{\operatorname{Log} z} = z$ for several test points.
\end{exercise}

\begin{exercise}
Show that the Taylor series for $e^z$ converges absolutely for all $z \in \mathbb{C}$.
\end{exercise}

\begin{exercise}
Verify the Cauchy-Riemann equations for $f(z) = e^z$ and $f(z) = \log z$.
\end{exercise}

\begin{exercise}
Prove that $e^{z_1 + z_2} = e^{z_1} \cdot e^{z_2}$ using the series definition.
\end{exercise}

\begin{exercise}
Find all solutions to $e^z = 1$ and explain why there are infinitely many.
\end{exercise}
